% ============================================================================
%  SECCIÓN 3.1: PROPUESTA DE LA SOLUCIÓN
% ============================================================================

\section{Propuesta de la soluci\'on}

De acuerdo con la gu\'ia del avance, en esta secci\'on se define la propuesta de
soluci\'on del proyecto: la identificaci\'on de los datos generales del equipo,
el nombre provisional y el tipo de aplicaci\'on, la descripci\'on de la soluci\'on
y el mapa de funcionalidades que permite visualizar de manera general los
m\'odulos que integrar\'an la aplicaci\'on m\'ovil.

\subsection{Datos del equipo}

El equipo de desarrollo est\'a compuesto por los siguientes integrantes:

\begin{table}[H]
  \centering
  \caption{Integrantes del equipo de desarrollo}
  \contenidoConFuente{%
    \begin{tablaAPA}
      \begin{tabular}{@{}p{5.6cm}p{5.0cm}p{5.0cm}@{}}
        \toprule
        \textbf{Integrante} & \textbf{Rol} & \textbf{Responsabilidad} \\
        \midrule
        Angel Gabriel Rojas Hinojosa & Desarrollador &
        Desarrollo de interfaces, integraci\'on de componentes y pruebas. \\
        \addlinespace
        Joaquin Alessandro Felipez Rojas & Desarrollador &
        Configuraci\'on del entorno, estructura del proyecto y control de
        versiones. \\
        \addlinespace
        Nicolas Sebastian Reguerin Meneses & Desarrollador &
        Desarrollo de pantallas, navegaci\'on y autenticaci\'on. \\
        \addlinespace
        David Willy Cruz Huanca & Desarrollador &
        Base de datos, servicios backend y configuraci\'on de Supabase. \\
        \bottomrule
      \end{tabular}
    \end{tablaAPA}%
  }{Elaboraci\'on propia.}
\end{table}

El representante del equipo es \textbf{Joaquin Alessandro Felipez Rojas},
 quien coordina las entregas y la comunicaci\'on con el docente de la
asignatura.

\subsection{Nombre provisional y tipo de aplicaci\'on}

La aplicaci\'on se denominar\'a provisionalmente \nombreApp{}, nombre con el que
se identifica el producto durante el desarrollo y que podr\'a ajustarse en
avances posteriores. En cuanto al tipo de soluci\'on, se propone el desarrollo
de una aplicaci\'on m\'ovil multiplataforma orientada a la
visualizaci\'on interactiva de escenarios deportivos a nivel nacional. La
aplicaci\'on m\'ovil podr\'a ejecutarse en los sistemas operativos iOS y Android
a partir de un \'unico c\'odigo base desarrollado con React Native y Expo, lo
que reduce los costos y los tiempos de desarrollo \citep{pressman2010}.

\subsection{Descripci\'on de la soluci\'on}

Se propone desarrollar una aplicaci\'on m\'ovil que permita a
la ciudadan\'ia consultar, localizar y acceder de forma centralizada a la
informaci\'on sobre los escenarios deportivos del pa\'is, incluyendo su
ubicaci\'on georreferenciada, sus caracter\'isticas, sus disciplinas, sus
horarios y sus eventos. La descripci\'on de la soluci\'on se organiza a
continuaci\'on respondiendo a las preguntas fundamentales de la propuesta:

\begin{enumerate}
  \item \textbf{\textquestiondown Qu\'e har\'a la aplicaci\'on?} La aplicaci\'on
    permitir\'a consultar un cat\'alogo de escenarios deportivos, visualizarlos
    en un mapa interactivo, revisar el detalle de cada recinto (disciplinas,
    horarios, servicios y eventos) y gestionar escenarios favoritos. Adem\'as,
    permitir\'a a los gestores publicar y actualizar la informaci\'on de sus
    escenarios y al administrador gestionar los usuarios y el cat\'alogo.

  \item \textbf{\textquestiondown Qu\'e problema ayudar\'a a resolver?}
    Contribuir\'a a resolver la dispersi\'on, la desactualizaci\'on y la falta
    de informaci\'on geoespacial sobre la infraestructura deportiva del pa\'is,
    que actualmente impide a la poblaci\'on conocer, localizar y acceder a la
    oferta deportiva nacional.

  \item \textbf{\textquestiondown Qui\'enes la utilizar\'an?} Ser\'a utilizada
    por los asistentes o aficionados deportivos, que consultan la oferta de
    escenarios; por los gestores de escenarios, que publican y actualizan la
    informaci\'on de sus recintos; y por los administradores del sistema, que
    gestionan los usuarios y validan la informaci\'on publicada, tal como se
    detalla en la secci\'on de identificaci\'on de usuarios.

  \item \textbf{\textquestiondown Qu\'e beneficios generar\'a?} Incrementar\'a
    la visibilidad de los escenarios y su aprovechamiento, promover\'a la
    pr\'actica deportiva y la actividad f\'isica, reducir\'a el tiempo y el
    esfuerzo que invierten los usuarios en la b\'usqueda de informaci\'on y
    apoyar\'a la planificaci\'on de las entidades administradoras.

  \item \textbf{\textquestiondown Por qu\'e una aplicaci\'on m\'ovil es adecuada
    para este problema?} Porque los tel\'efonos inteligentes constituyen el medio
    de acceso a la informaci\'on m\'as extendido entre la poblaci\'on, permiten
    aprovechar la geolocalizaci\'on para encontrar los escenarios m\'as cercanos
    al usuario y facilitan la consulta de la informaci\'on en el lugar y en el
    momento en que se necesita \citep{pressman2010, longley2015}.
\end{enumerate}

\subsection{Mapa de funcionalidades}

El mapa de funcionalidades presenta la visi\'on general del sistema, agrupando
los m\'odulos funcionales que componen la aplicaci\'on \nombreApp{} y que se
desarrollan con mayor detalle en las secciones de funcionalidades preliminares
y de definici\'on del MVP. En la figura se muestra la estructura general de
m\'odulos de la aplicaci\'on m\'ovil.

\begin{figure}[H]
\centering
  \figuraAPA{fig:mapa-funcionalidades}{Mapa de funcionalidades de la aplicaci\'on \nombreApp{}}{%
  \begin{tikzpicture}[
      nucleo/.style={rectangle, draw=black!80, thick, rounded corners=3pt,
                     align=center, fill=orange!15, text width=3.6cm,
                     minimum height=1.3cm, font=\small\bfseries},
      modulo/.style={rectangle, draw=black!60, rounded corners=2pt,
                     align=center, fill=blue!8, font=\scriptsize,
                     text width=2.9cm, minimum height=0.85cm},
      flecha/.style={-{Stealth[length=2mm]}, thick, draw=black!50}
    ]
    \node[nucleo] (nucleo) at (0,0) {APLICACI\'ON \nombreApp{}\\ (m\'ovil)};
    \node[modulo] (auth) at (0,2.2) {Registro e inicio de sesi\'on};
    \node[modulo] (cat) at (-4.2,0.9) {Cat\'alogo de escenarios};
    \node[modulo] (bus) at (4.2,0.9) {B\'usqueda y filtros};
    \node[modulo] (mapa) at (-4.2,-1.1) {Mapa interactivo};
    \node[modulo] (det) at (4.2,-1.1) {Detalle del escenario};
    \node[modulo] (fav) at (-4.2,-2.9) {Favoritos};
    \node[modulo] (gest) at (0,-2.9) {Gesti\'on de contenido};
    \node[modulo] (notif) at (4.2,-2.9) {Notificaciones};

    \draw[flecha] (auth) -- (nucleo);
    \draw[flecha] (cat) -- (nucleo);
    \draw[flecha] (bus) -- (nucleo);
    \draw[flecha] (mapa) -- (nucleo);
    \draw[flecha] (det) -- (nucleo);
    \draw[flecha] (fav) -- (nucleo);
    \draw[flecha] (gest) -- (nucleo);
    \draw[flecha] (notif) -- (nucleo);
  \end{tikzpicture}%
  }{Elaboraci\'on propia en base a las funcionalidades preliminares del sistema.}
\end{figure}

Los m\'odulos presentados se desarrollan con mayor detalle en las secciones de
funcionalidades preliminares y de definici\'on del MVP, donde se establece
cu\'ales de ellos forman parte del producto m\'inimo viable y cu\'ales se
incorporar\'an como mejoras en avances posteriores.
